\section{Metodologie e Dati}\label{sec:methods}

\subsection{I dati}
I dati sono stati estratti da un dataset storico completo delle Olimpiadi
moderne presente su Kaggle, integrato dal 2017 in poi dai dati presenti su
Wikipedia (1924--2024), con informazioni su:
\begin{itemize}
	\item Paese partecipante (codice NOC --- National Olympic Committee)
	\item Anno dell'evento olimpico
	\item Status di ospitalità (binario)
	\item Numero totale di medaglie vinte (oro, argento, bronzo sommate)
\end{itemize}
Le medaglie vengono contate per evento, vale a dire che gli sport di squadra
possono fornirne una per team, piuttosto che una per ogni suo membro.

\subsection{Il metodo}
Il test chi-quadrato di Pearson è stato utilizzato per testare l'indipendenza
tra due variabili categoriche: ospitalità (sì/no) e numero di medaglie.\\
Il test è stato applicato individualmente a ognuno dei 18 Paesi che abbiano
ospitato almeno un'edizione nel periodo selezionato. Per i paesi che non
soddisfacevano l'assunzione $\textit{frequenze\_attese} \ge 5$ (Grecia e
Messico), i risultati sono riportati con un avviso di cautela (*), in quanto non
statisticamente affidabili secondo gli standard di Cochran (1954).\\
\\
Inoltre, è stato applicato in maniera aggregata (per somma) a tutti i Paesi così
selezionati e che soddisfacessero il teorema di Cochran (16 in totale).\\
Il risultato aggregato fornirà una valutazione generale dell'effetto host,
mentre quelli individuali delle singole Nazioni permetteranno di comprendere
quali, di preciso, ne abbiano beneficiato.\\
\\
Seguono le formule usate, prima per ogni Paese, poi per tutti in maniera aggregata:\\
\begin{itemize}
	\item $O_i$ = frequenza osservata, di medaglie, rispettivamente in casa e fuori
	\item $E_i$ = frequenza attesa sotto ipotesi nulla (in base alla media storica di medaglie)
	\item Frequenze attese:
\end{itemize}
\[
	E_{ospita}=\text{Medaglie\_Totali} \times \frac{n_{ospita}}{n_{totale}}
\]
\[
	E_{\text {non\_ospita}}=\text{Medaglie\_Totali} \times \frac{n_{non\_ospita}}{n_{totale}}
\]
\begin{itemize}
	\item Chi quadro:
\end{itemize}
\[
	\chi^2= \sum_{i=1}^{2} \frac{(O_i - E_i)^2}{E_{i}}
\]
